\scanpage{36}
% archive sheet number: 35
le $\overline{\mathrm m}$ argument montre que pour tte $K$-extension finie $L$ de $\overline K/K$, $\varphi$ induit
\[
\varphi_L:X(L)\times M_L\longrightarrow\mathbb R\quad,
\]
ces applications $\varphi_L$ ont une propriété de transitivité évidente
\[
\varphi_{L^{\deleted{\unsure{\alpha|*|'}}}}(P,v|L)\simeq\varphi_{L'}(P,v)
\qquad\left|\begin{array}{l}P\in X(L)\\ v\in M_{L'}\end{array}\right.
\]
et la donnée de $\varphi$ équivaut au système des $\varphi_L$ satisfaisant la condition de compatibilité précédente. On \deleted{dit} désigne par $\widehat{\Phi}(X)$ [ou, en cas d'\unsure{ambiguité|ambiguïté}, $\widehat{\Phi}(X/K,M_K)$] l'ens de ces $\deleted{\unsure{\Phi}}\varphi$. Noter que $\widehat{\Phi}(X)$ est déterminé à bij. unique près, indépendamment du choix de la clôture alg. $\overline K$.

On a, pour une extension \inserted{finie} $K'$ de $K$, une application naturelle
\[
\widehat{\Phi}(X/\deleted{\unsure{X}}K)\longrightarrow\widehat{\Phi}(X_{K'}/K')
\]
\deleted{et pour $X$ fixé,}, qui est \uline{bijective}
